\clearpage
\VerdierPhysicalPageAnchor{0235}%
\begingroup\small\setlength{\parskip}{0pt}\setlength{\emergencystretch}{1.5em}
\setlength{\abovedisplayskip}{4pt plus 2pt minus 2pt}%
\setlength{\belowdisplayskip}{4pt plus 2pt minus 2pt}%

\VerdierSemanticAnchor{ch4d-4-2-2-cont}
\noindent 其中 $(p,q)$ 与 $(p',q')$ 是
$\mathcal{F}\ell_1(\widetilde{\mathbb Z})$ 的元素，且 $p\leq p'$、
$q\leq q'$；上式表示 $\operatorname{comp}(\mathcal A)$ 中典范态射

\VerdierSemanticAnchor{ch4d-4-2-2-3}
\begin{equation}
\psi\bigl((p,q),(p',q')\bigr)\colon
X(p,q)\longrightarrow X(p',q').
\tag{4.2.2.3}
\end{equation}
在 $D(\mathcal A)$ 中的像。这样就确定了一个从
$\mathcal{F}\ell_1(\widetilde{\mathbb Z})$ 到 $D(\mathcal A)$ 的函子。
对 $\mathcal{F}\ell_2(\widetilde{\mathbb Z})$ 的每个元素 $(p,q,r)$，令

\VerdierSemanticAnchor{ch4d-4-2-2-4}
\begin{equation}
S(p,q,r)\colon
0\longrightarrow X(p,q)\longrightarrow X(p,r)
\longrightarrow X(q,r)\longrightarrow0
\tag{4.2.2.4}
\end{equation}
为 $\operatorname{comp}(\mathcal A)$ 中的正合列，其中的态射都是由 $X$ 的
滤过所导出的典范态射。最后，对
$\mathcal{F}\ell_1(\widetilde{\mathbb Z})$ 的每个元素 $(p,q)$，令

\VerdierSemanticAnchor{ch4d-4-2-2-5}
\begin{equation}
J(p,q)\colon\mathcal X(p,q)\longrightarrow X(p,q)
\tag{4.2.2.5}
\end{equation}
为 $\operatorname{comp}(\mathcal A)$ 中态射

\VerdierSemanticAnchor{ch4d-4-2-2-6}
\begin{equation}
j(p,q)\colon\mathcal X(p,q)\longrightarrow X(p,q),
\tag{4.2.2.6}
\end{equation}
在 $D(\mathcal A)$ 中的像；其分量为：

\VerdierSemanticAnchor{ch4d-4-2-2-7}
\begin{equation}
\left\{
\begin{aligned}
j(p,q)^i\colon\mathcal X(p,q)^i
 &=X(q)^i\oplus X(p)^{i+1}
   \longrightarrow X(q)^i/X(p)^i=X(p,q)^i,\\
j(p,q)^i&=(\pi^i,0),\\
&\text{其中 $\pi^i$ 是典范满态射：}\\
\pi^i\colon X(q)^i&\longrightarrow X(q)^i/X(p)^i.
\end{aligned}
\right.
\tag{4.2.2.7}
\end{equation}

\VerdierSemanticAnchor{ch4d-4-2-3}
\medskip\noindent\textbf{命题 4.2.3.}\quad
\textit{对 $\mathcal{F}\ell_2(\widetilde{\mathbb Z})$ 的每个元素 $(p,q,r)$，记}
\[
\delta(p,q,r)\colon X(q,r)\longrightarrow X(p,q)[1]
\]
\noindent\textit{为 $D(\mathcal A)$ 中的态射}

\VerdierSemanticAnchor{ch4d-4-2-3-1}
\begin{equation}
\delta(p,q,r)=\delta\bigl(S(p,q,r)\bigr),
\tag{4.2.3.1}
\end{equation}
\noindent\textit{其中 $\delta\bigl(S(p,q,r)\bigr)$ 是典范 $\delta$-函子
$\operatorname{comp}(\mathcal A)\longrightarrow D(\mathcal A)$ (1.3.2)
关联给正合列 $S(p,q,r)$ (4.2.2.4) 的 $D(\mathcal A)$ 中的态射。}

\VerdierSemanticAnchor{ch4d-4-2-3-ab}
\noindent\textit{a) 对象 $X(p,q)$ (4.2.2.1)、态射
$\varphi\bigl((p,q),(p',q')\bigr)$ (4.2.2.2) 和态射
$\delta(p,q,r)$ (4.2.3.1) 定义了一个取值于 $D(\mathcal A)$ 的谱对象。}

\noindent\textit{b) 态射 $J(p,q)$ (4.2.2.5) 定义了从谱对象
$\bigl(\mathcal X(p,q),\delta(p,q,r)\bigr)$ 到谱对象
$\bigl(X(p,q),\delta(p,q,r)\bigr)$ 的同构。}

\VerdierSemanticAnchor{ch4d-4-2-3-proof}
a) 由典范 $\delta$-函子
$\operatorname{comp}(\mathcal A)\longrightarrow D(\mathcal A)$ 的性质
(1.3.1) 立即得出。为证明 b)，读者可参阅 (1.3.2)。
\endgroup
